\subsection{Radial multiplication preserving the divergence sign}

\begin{proposition}[Radial multiplication of a homogeneous potential]
\label{prop:cubic-radial-lift}
Let \(h\) be smooth and homogeneous of degree \(d>0\) on a conical
open subset of \(\R^N\setminus\{0\}\). If
\[
 h\cL(h)\ge0,\qquad h\Delta h\ge0,\qquad N\ge2d+2,
\]
then \(g=\sqrt R\,h\) satisfies \(g\cL(g)\ge0\), strictly where
\(h\ne0\). If \(\nabla h\ne0\) throughout the domain, then
\(\nabla g\ne0\) there as well.
\end{proposition}

\begin{proof}
The radial master identity, \cref{prop:radial-master-identity} with
\(\alpha=1/2\), reads
\begin{align}
 \cL(\sqrt R\,h)=R^{3/2}\Biggl[\cL(h)
 &+(2d+1)\frac{h^2\Delta h}{R}
 +(N-2d-2)\frac{h|\nabla h|^2}{R}\notag\\
 &+(2d+1)(d+N-1)\frac{h^3}{R^2}\Biggr].
 \label{eq:cubic-radial-lift}
\end{align}
Multiplication by \(g=\sqrt R\,h\) makes every term nonnegative;
the last one is positive when \(h\ne0\). The same master identity
gives
\[
 |\nabla g|^2=R|\nabla h|^2+(2d+1)h^2,
\]
proving the final assertion.
\end{proof}
